\documentclass[reqno,11pt]{amsart}
\usepackage{math327-lecture}
\usepackage{needspace}
\newcommand{\C}{\mathbb{C}}
\lectureinfo{25}{Nov. 16, 2026}{Rings, Units, Ring Homomorphisms, and Polynomial Rings}
\artinsections{11.1--11.3}

% Typeset from handwritten Lecture 27, Nov. 17, 2025, pp. 1--13.
\begin{document}
\raggedbottom
\makelecturetitle



\section{Rings}

We now turn to ring theory. A ring is a set equipped with two
binary operations, an addition and a multiplication, which interact
in the way familiar from the integers.

\begin{define}[Ring]
A \emph{ring} $(R,+,\cdot,0,1)$ is a set $R$ together with two
binary operations
\[
  +:R\times R\to R,\quad (a,b)\mapsto a+b,
  \qquad
  \cdot:R\times R\to R,\quad (a,b)\mapsto a\cdot b,
\]
and two elements $0,1\in R$, such that for all $a,b,c\in R$:
\begin{itemize}
  \item $(R,+,0)$ is an abelian group; that is, $+$ is associative,
    $a+b=b+a$, $a+0=0+a=a$, and each $a$ has an additive inverse
    $-a\in R$ with $a+(-a)=0$;
  \item multiplication is associative and commutative:
    $a\cdot(b\cdot c)=(a\cdot b)\cdot c$ and $a\cdot b=b\cdot a$;
  \item $1$ is a multiplicative identity: $a\cdot1=1\cdot a=a$;
  \item multiplication distributes over addition:
    $a\cdot(b+c)=a\cdot b+a\cdot c$.
\end{itemize}
\end{define}

\begin{remark}
In general $(R,\cdot,1)$ is \emph{not} a group: elements need not
have multiplicative inverses. The distributive law is what ties the
two operations together. In this course rings are commutative and
have a multiplicative identity $1$.
\end{remark}

\begin{example}[The integers]
$\Z=(\Z,+,\cdot,0,1)$ is a ring. The additive group $(\Z,+,0)$ is
the usual group we have been dealing with, but $\Z$ also has a
multiplication: if $m,n\in\Z$, then $m\cdot n\in\Z$.
\end{example}

\section{Units}

In general an integer $m$ does not have an ``inverse'' under
multiplication. However,
\[
  1\cdot1=1,\qquad (-1)\cdot(-1)=1,
\]
so $+1$ and $-1$ do have multiplicative inverses.

\begin{define}[Group of units]
Let $R$ be a ring. The \emph{group of units} of $R$ is
\[
  R^\times=\{a\in R:\text{there exists } b\in R\text{ with }
  ab=ba=1\}.
\]
\end{define}

\begin{proposition}\label{prop:units-group}
$R^\times\subseteq R$ is a group under multiplication.
\end{proposition}
\begin{proof}
\begin{enumerate}[label=(\arabic*)]
  \item $1\in R^\times$, since $1\cdot1=1$.
  \item $R^\times$ is closed under multiplication: if
    $a_1,a_2\in R^\times$, choose $b_1,b_2$ with $a_1b_1=1$ and
    $a_2b_2=1$. Then
    \[
      (a_1a_2)(b_1b_2)=(a_1b_1)(a_2b_2)=1\cdot1=1,
    \]
    so $a_1a_2\in R^\times$.
  \item If $a\in R^\times$, then $a$ has an inverse by definition,
    and that inverse lies in $R^\times$ as well.
\end{enumerate}
Associativity is inherited from $R$.
\end{proof}

Thus a ring $R$ gives rise to two abelian groups:
\[
  (R,+,0)\qquad\text{and}\qquad(R^\times,\cdot,1).
\]

\begin{example}
For $R=\Z$, these are $(\Z,+,0)$ and
\[
  (\{\pm1\},\cdot,1)\cong\Z/2\Z.
\]
\end{example}

\section{First examples: Gaussian integers, the zero ring, subrings}

\begin{example}[Gaussian integers]
Inside $\C$, where $i^2=-1$, consider
\[
  \Z[i]=\{a+ib:a,b\in\Z\}\subseteq\C.
\]
To make this a ring we need $+$, $\cdot$, $0$, and $1$. Define
\begin{align*}
  (a+ib)+(\tilde a+i\tilde b)&=(a+\tilde a)+i(b+\tilde b),\\
  (a+ib)\cdot(\tilde a+i\tilde b)
  &=a\tilde a+ib\tilde a+i^2b\tilde b+ia\tilde b
   =(a\tilde a-b\tilde b)+i(b\tilde a+a\tilde b),
\end{align*}
using $i^2=-1$, and take $0=0+i\cdot0$ and $1=1+i\cdot0$.
\end{example}

\begin{exercise}
Check that $\Z[i]$, the \emph{ring of Gaussian integers}, with the
operations $+$ and $\cdot$ above and with $0=0+i\cdot0$,
$1=1+i\cdot0$, is a ring.
\end{exercise}

What happens in a ring $(R,+,\cdot,0,1)$ if $0=1$?

\begin{proposition}[The zero ring]\label{prop:zero-ring}
If $0=1$ in a ring $R$, then $R=\{0\}$.
\end{proposition}
\begin{proof}
Pick $a\in R$. Since $0=1-1$, distributivity gives
\[
  0\cdot a=(1-1)a=a-a=0.
\]
Hence
\[
  a=1\cdot a=0\cdot a=0.
\]
Conversely, $\{0\}$ is a ring with $0+0=0$ and $0\cdot0=0$
(here $1=0$).
\end{proof}

\begin{define}[Subring]
A subset $S$ of a ring $R$ is a \emph{subring} if it is closed under
$+$ and $\cdot$ and contains $0$ and $1$.
\end{define}
% NOTE: the notes say only "closed under +, . and contains 0, 1";
% closure under additive inverses (needed for (S,+,0) to be a group)
% is left implicit there.

\begin{example}
$\C$ is a ring under the usual addition and multiplication of
complex numbers, and $\Z[i]\subseteq\C$ is a subring.
\end{example}

\section{Ring homomorphisms}

Before more examples, let us introduce the notion of a ring
homomorphism.

\begin{define}[Ring homomorphism]
Let $R=(R,+_R,\cdot_R,0_R,1_R)$ and $S=(S,+_S,\cdot_S,0_S,1_S)$ be
two rings. A \emph{ring homomorphism} is a map $\varphi:R\to S$ such
that for all $a,b\in R$,
\begin{align*}
  \varphi(a+_Rb)&=\varphi(a)+_S\varphi(b),\\
  \varphi(a\cdot_Rb)&=\varphi(a)\cdot_S\varphi(b),\\
  \varphi(1_R)&=1_S.
\end{align*}
\end{define}

\begin{remark}
The condition $\varphi(0_R)=0_S$ follows from the first condition,
exactly as for group homomorphisms. The identity
$\varphi(1_R)=1_S$, however, must be imposed separately.
\end{remark}

\begin{remark}[Units go to units]
If $\varphi:R\to S$ is a ring homomorphism, then
$\varphi(R^\times)\subseteq S^\times$. Indeed, let $a\in R^\times$
and choose $b\in R^\times$ with $ab=1$. Then
\[
  1=\varphi(1)=\varphi(ab)=\varphi(a)\varphi(b),
\]
so $\varphi(a)\in S^\times$ and $\varphi(a)^{-1}=\varphi(a^{-1})$.
\end{remark}

\begin{example}[Homomorphisms out of $\Z$]\label{ex:Z-to-R}
What is a ring homomorphism $\Z\to R$? It must send $1\mapsto1_R$,
and then additivity forces
\begin{align*}
  2=1+1&\longmapsto 1_R+1_R,\\
  3&\longmapsto 1_R+1_R+1_R,\\
  -1&\longmapsto -1_R,\\
  -2&\longmapsto (-1_R)+(-1_R),
\end{align*}
and so on. Here $-1_R$ means the additive inverse of $1_R$ in
$(R,+,0)$. So there is \emph{at most one} ring homomorphism
$\Z\to R$, and one checks that the map just described actually
satisfies the conditions of a homomorphism. Hence there is a
unique ring homomorphism $\Z\to R$.
\end{example}

\section{Polynomial rings}

Our main example is the following.

\begin{define}[Polynomial ring]
Let $R$ be a ring. The \emph{ring of polynomials over $R$} is
\[
  R[x]=\{f=a_0+a_1x+\cdots+a_nx^n:a_i\in R,\ n\geq0\}.
\]
\end{define}

\begin{define}[Degree, leading coefficient, monic]
Let $f=\sum_i a_ix^i\in R[x]$ be nonzero.
\begin{itemize}
  \item The \emph{degree} $\deg(f)$ is the largest $n$ with
    $a_n\neq0$.
  \item The \emph{leading} (or \emph{highest}) \emph{coefficient}
    of $f$ is $a_n$ for this largest $n$.
  \item $f$ is \emph{monic} if its leading coefficient is $1$, that is,
    \[
      f=x^n+a_{n-1}x^{n-1}+\cdots+a_0.
    \]
\end{itemize}
\end{define}

\begin{proposition}
$R[x]$ is a ring under the following operations. For
$f=\sum_ia_ix^i$ and $g=\sum_jb_jx^j$,
\[
  f+g=\sum_i(a_i+b_i)x^i,\qquad
  f\cdot g=\sum_k\Bigl(\sum_{i+j=k}a_ib_j\Bigr)x^k,
\]
with zero element the constant polynomial $0$ and identity the
constant polynomial $1$.
\end{proposition}

Most of the interesting facts about $R[x]$ are deduced from division
with remainder. Recall that in $\Z$, given $m$ and $n>0$, we can
write $m=an+r$ with $r=0$ or $0<r<n$. In $R[x]$ this can fail:
in $\Z[x]$ we cannot divide $7x+1$ by $2x$ with a remainder of
smaller degree, because $7/2\notin\Z$. The fix is to divide only by
polynomials with a suitable leading coefficient.

\begin{theorem}[Division with remainder]\label{thm:division}
Let $f(x)\in R[x]$ be monic. Then for any $g(x)\in R[x]$ there exist
unique $q(x),r(x)\in R[x]$ such that
\[
  g=q\cdot f+r,
\]
where $r=0$ or $\deg(r)<\deg(f)$.
\end{theorem}

\begin{remark}
If we replace ``$f$ monic'' by ``the leading coefficient of $f$ is a
unit'', that is, lies in $R^\times$, Theorem~\ref{thm:division} still
works.
\end{remark}

In Math 30-1 you work with $\R[x]$. There leading coefficients lie
in $\R\setminus\{0\}$, but every such element is a unit, so division
with remainder is always possible. In other words, $\R$ is a field.

\begin{define}[Field]
A ring $R$ is a \emph{field} if $R^\times=R\setminus\{0\}$, that is,
if every nonzero element of $R$ has a multiplicative inverse.
\end{define}

\begin{question}
What is a ring homomorphism $R[x]\to S$ which extends a given
homomorphism $\varphi:R\to S$? For instance, there is a unique
homomorphism $\Z\to R$ by Example~\ref{ex:Z-to-R}; what are the
homomorphisms $\Z[x]\to\Z$?
\end{question}

\end{document}
